% from ega4-2-fr.tex
\oldpage[IV-2]{67}

Or, il existe un corps $K$, extension commune de $k'$ et de $k(s'')$, de sorte
que si l'on pose $P=\operatorname{Spec}(K)$, le diagramme
\[
\xymatrix@C=3.5em@R=3em{
S'\ar[d] & P\ar[l]\ar[d]\\
S & T\ar[l]
}
\]
soit commutatif~; si l'on pose $Q=X\otimes_k K$, le diagramme
\[
\xymatrix@C=3.5em@R=3em{
X'\ar[d]_p & Q\ar[l]_q\ar[d]^r\\
X & U\ar[l]^w
}
\]
est donc lui aussi commutatif, et les deux composantes irréductibles
$r^{-1}(U_1)$ et $r^{-1}(U_2)$ de $Q$ (4.4.5) contiennent par suite
$q^{-1}(Z')$. Or, en vertu de (4.4.5), les images
$q(r^{-1}(U_1))$ et $q(r^{-1}(U_2))$ sont des composantes irréductibles de $X'$
contenant $Z'$~; elles sont donc toutes deux égales à $X'_0$, donc
$r^{-1}(U_1)=r^{-1}(U_2)$ en vertu de (4.4.5), et finalement $U_1=U_2$.

Comme $p$ est surjectif et que $X'_0$ domine une composante irréductible de
$X$, $X_0=p(X'_0)$ est une composante irréductible de $X$ contenant $Z$, et
comme $p^{-1}(X_0)$ est une composante irréductible de $X'$, $X_0$ est
géométriquement irréductible.

Pour prouver la dernière assertion, soit $X_1$ une seconde composante
irréductible de $X$ rencontrant $Z$. Prenant $z\in Z\cap X_1$, on voit qu'il
existe une composante irréductible de $X'$ rencontrant $Z'$ et dominant $X_1$
(2.3.5)~; mais comme $X'_0$ est par hypothèse la seule composante irréductible
de $X'$ rencontrant $Z'$, on a nécessairement $X_1=X_0$.
\end{proof}

\begin{remark}[4.5.20]
\label{IV.2.4.5.20-fr}
Le fait que $Z'$ soit contenu dans une seule composante irréductible $X'_0$ de
$X'$ \emph{n'entraîne pas} que $X_0=p(X'_0)$ soit la seule composante
irréductible de $X$ contenant $Z$. Prenons en effet $k=\mathbf R$,
$k'=\mathbf C$, et pour $X$ le $k$-schéma obtenu par recollement\footnote{La
technique de «~recollement~» des préschémas sera développée en détail au
chap.~V~; pour le cas considéré ici (où il s'agit de courbes algébriques), on
peut consulter [38], p.~68--71.} en un point (non générique) de
$X_1=\operatorname{Spec}(\mathbf R[S,T]/(S^2+T^2+1))$ et de
$X_2=\operatorname{Spec}(\mathbf C[T])$ (les anneaux locaux de ces deux courbes
en chacun de leurs points fermés ayant même corps résiduel isomorphe à
$\mathbf C$)~; $X_1$ et $X_2$ s'identifient aux deux composantes irréductibles
de $X$~; si $x$ est leur point commun, $p^{-1}(x)$ est formé de deux points
distincts $y'$, $z'$, $p^{-1}(X_1)$ est irréductible, et $p^{-1}(X_2)$ est somme
de deux composantes irréductibles $Y'$, $Z'$ de $X'$ telles que $y'\in Y'$ et
$z'\in Z'$, de sorte que $p^{-1}(x)$ n'est contenu que dans une seule
composante irréductible de $X'$.
\end{remark}

\begin{proposition}[4.5.21]
\label{IV.2.4.5.21-fr}
Soit $k$ un corps séparablement clos, et soit $k'$ une clôture algébrique de
$k$. Pour tout $k$-préschéma $X$, la projection canonique
$X\otimes_k k'\to X$ est un homéomorphisme universel.

